\chapter{Annotations}
\label{ch:annotations}

Annotations are rosetta's \textbf{opt-in enrichment layer}: small declarative
tags attached to your fields and methods that backends turn into docstrings,
validation and user interface. They are never required --- reflection alone
binds the whole class --- and they never change the C++ behaviour of your types.

They buy three things:

\begin{enumerate}
  \item \textbf{Documentation} --- \code{doc} becomes docstrings, reference
    pages, tooltips.
  \item \textbf{Validation} --- \code{range} rejects bad assignments at the
    binding boundary.
  \item \textbf{User interface} --- \code{label}, \code{button},
    \code{combobox} and the \code{widget::*} hints drive the generated Qt, QML
    and Dear ImGui inspectors.
\end{enumerate}

\section{Two spellings, one result}

Every annotation can be written \textbf{inline} (as a P2996/P3394 attribute in
the header) or \textbf{out of line} (as a JSON side-car, leaving the header
plain C++). The two forms produce identical bindings and can be mixed on the
same class.

\begin{lstlisting}[style=cpp]
#include <rosetta/annotations.h>

class Algo {
public:
    [[= rosetta::doc{"Max solver iterations"}, = rosetta::range{0, 200},
       = rosetta::widget::slider, = rosetta::label{"Max iter"} ]]
    int maxIter{100};

    [[= rosetta::doc{"Run the solver"}, = rosetta::button{"Run"} ]]
    double run();
};
\end{lstlisting}

\ldots or the same thing with a clean header and a side-car:

\begin{lstlisting}[style=json]
{
  "maxIter": { "doc": "Max solver iterations", "range": [0, 200],
               "widget": "slider", "label": "Max iter" },
  "run":     { "doc": "Run the solver", "button": "Run" }
}
\end{lstlisting}

Annotations are read \emph{once}, on the generation host, and merged (inline
first, then JSON) into the same annotation list. Backends cannot tell the two
sources apart.

\section{The annotation set}

\begin{center}
\small
\begin{longtable}{@{}lL{2.0cm}L{5.7cm}L{4.9cm}@{}}
\toprule
\textbf{Kind} & \textbf{Applies to} & \textbf{Spellings} & \textbf{Meaning} \\
\midrule
\endhead
\code{doc} & fields, methods, classes$^*$ &
  \code{= rosetta::doc\{"\ldots"\}} \newline \code{"doc": "\ldots"} &
  Description text. Docstrings, reference pages, OpenAPI descriptions, hover
  tooltips. \\
\addlinespace
\code{range} & numeric fields &
  \code{= rosetta::range\{lo, hi\}} \newline \code{"range": [lo, hi]} &
  Value bounds. Compiled backends emit a \textbf{validating setter}; UI
  inspectors render a clamping slider; doc backends print the bound. Scientific
  notation is fine (\code{[1e-10, 1e-6]}). \\
\addlinespace
\code{readonly} & fields &
  \code{= rosetta::readonly} \newline \code{"readonly": true} &
  Getter-only property; writes are rejected per backend; UI editors grey out. \\
\addlinespace
\code{combobox} & string fields &
  \code{= rosetta::combobox\{\{"a","b"\}\}} \newline
  \code{"combobox": ["a","b"]} &
  Allowed choices (maximum 16). UI renders a drop-down. \\
\addlinespace
\code{label} & fields, methods &
  \code{= rosetta::label\{"\ldots"\}} \newline \code{"label": "\ldots"} &
  Display-name override in UI backends. The C++ identifier stays the
  lookup/binding key --- only the visible text changes. \\
\addlinespace
\code{button} & methods &
  \code{= rosetta::button\{"Run"\}} \newline \code{"button": "Run"} &
  Caption of the call button a UI backend renders for the method. \\
\addlinespace
\code{widget::*} & fields &
  \code{= rosetta::widget::slider} \newline \code{"widget": "slider"} &
  Editor hint --- see below. \\
\bottomrule
\end{longtable}
\end{center}

\noindent$^*$ Class-level and free-function \code{doc} come from the
\textbf{manifest} (a \field{doc} key on the entry), since free functions carry
no in-source annotations.

\subsection{Widget hints}

Hints pick the editor when more than one would fit the field's type. They are
consumed by the three UI inspector backends and ignored everywhere else.

\begin{center}
\small
\begin{tabular}{@{}llL{7.0cm}@{}}
\toprule
\textbf{Hint} & \textbf{Field type} & \textbf{Renders as} \\
\midrule
\code{spin}      & number & spin box / typed input \\
\code{slider}    & number + \code{range} & slider with the range as bounds
  (clamps by construction) \\
\code{textfield} & number & validated free-text input --- the right editor for
  values a slider cannot hit precisely (an eps of \code{1e-7}) \\
\code{checkbox}  & number (0/1), bool & checkbox \\
\code{color}     & string \code{"\#rrggbb"} & colour picker (ImGui
  \code{ColorEdit3}, Qt swatch + \code{QColorDialog}, QML \code{ColorDialog});
  edits write back as hex \\
\code{multiline} & string & multi-line text area \\
\code{radio}     & string + \code{combobox} & the choices as a radio-button
  group instead of a drop-down \\
\code{file}      & string path & text entry + browse button opening the platform
  file dialog \\
\bottomrule
\end{tabular}
\end{center}

\section{Who consumes what}

\begin{description}[leftmargin=0cm,style=nextline,font=\normalfont\itshape]
  \item[Scripting backends] (\lang{python}, \lang{nanobind}, \lang{node},
    \lang{wasm}, \lang{lua}, \lang{julia}, plus \lang{csharp} / \lang{java}):
    \code{doc} becomes docstrings where the framework has them; \code{range}
    becomes validating setters; \code{readonly} becomes getter-only properties.
    \code{label}, \code{button} and \code{widget::*} are ignored --- they are UI
    concepts.
  \item[UI inspectors] (\lang{qt}, \lang{qml}, \lang{imgui}): everything.
    \code{doc} as tooltips, \code{range} as slider bounds or commit validation,
    \code{readonly} as disabled editors, \code{combobox} as drop-downs, plus
    \code{label}, \code{button} and \code{widget::*}.
  \item[Document backends] (\lang{markdown}, \lang{html}, \lang{openapi},
    \lang{typescript}): \code{doc} and the constraints (\code{range},
    \code{readonly}, \code{combobox}) are rendered into the reference text or
    schema.
\end{description}

An annotation a backend has no use for is simply ignored --- never an error.

\section{Out-of-line annotations}
\label{sec:outofline}

The side-car keeps your header plain C++. That matters for three reasons, and
the second is the one that usually decides it:

\begin{center}
\small
\begin{tabular}{@{}L{6.6cm}cc@{}}
\toprule
 & \textbf{Inline} & \textbf{Side-car} \\
\midrule
Header stays plain C++ & \no & \yes \\
Generated binding builds on an off-the-shelf compiler & \no & \yes \\
Works on third-party headers you cannot edit & \no & \yes \\
Expressiveness & full set & full set \\
\bottomrule
\end{tabular}
\end{center}

The middle row is the load-bearing one. Inline annotations pull
\code{<experimental/meta>} into your header, and the generated binding
\code{\#include}s that header --- so the ``reflection-free by construction''
promise (Section~\ref{sec:reflection-free}) is lost for every target.

\subsection{Wiring it up}

One manifest field per class:

\begin{lstlisting}[style=json]
"classes": [
  { "name": "Algo", "header": "Algo.h", "annotations": "Algo.ann.json" }
]
\end{lstlisting}

The path is resolved relative to the manifest. The side-car's keys are member
identifiers --- fields \emph{or} methods --- and each value object may carry any
of the annotation kinds:

\begin{lstlisting}[style=json]
{
  "title": { "doc": "The widget title", "label": "Title" },
  "count": { "doc": "Visible items", "range": [0, 100], "widget": "slider" },
  "id":    { "readonly": true },
  "mode":  { "combobox": ["fast", "slow"] },
  "run":   { "doc": "Run the thing", "button": "Run" }
}
\end{lstlisting}

The bound header never mentions rosetta:

\begin{lstlisting}[style=cpp]
#pragma once
#include <string>
struct Widget {
    std::string title;
    int         count = 0;
    int         id    = 0;
    std::string mode;
};
\end{lstlisting}

That is the whole surface. Run \code{rosetta\_gen}, build, and every backend
sees the annotations.

\subsection{What it does under the hood}

\code{rosetta\_gen} bakes the JSON bytes into \code{bindings.h}, right after the
class header is included so \code{T} is complete:

\begin{lstlisting}[style=cpp]
template <> inline constexpr auto rosetta::detail::ann_storage<Widget> =
    std::to_array<char>({ char(0x7b), /* baked JSON bytes */ '\0' });
template <> inline constexpr std::string_view rosetta::ann_json_source<Widget> =
    std::string_view{ rosetta::detail::ann_storage<Widget>.data(),
                      rosetta::detail::ann_storage<Widget>.size() - 1 };
static_assert(rosetta::detail::ann_keys_error<Widget>().empty(),
              rosetta::detail::ann_keys_error<Widget>());
\end{lstlisting}

A small \code{consteval} parser reads it at \code{walk()} time and concatenates
the entries with the member's inline annotations. Nothing downstream changes.

Loading the JSON at \emph{run time} could not work: annotations are non-type
template parameters consumed during constant evaluation, so the data must exist
then. Baking keeps everything in the one constant-evaluated pass --- and keeps
the \code{static\_assert} below possible.

\subsection{The staleness guard}

That \code{static\_assert} checks every JSON key names a real member of the
type. A renamed or misspelt field \textbf{fails the build}, naming the offending
key, instead of silently dropping its metadata.

\begin{warn}
Side-car edits need a \textbf{regeneration}, not just a rebuild --- the JSON is
baked when \code{rosetta\_gen} runs. Re-run \code{rosetta\_gen} and the
generator, not only the binding build.
\end{warn}

\subsection{Attaching the JSON yourself}

The manifest's \field{annotations} field is sugar for specialising a
customisation point. You can do it directly, from anywhere --- another file, a
translation unit, after the class:

\begin{lstlisting}[style=cpp]
#include "widget.h"            // the untouched class
#include <rosetta/annotate.h>  // for the customization point

template <>
constexpr std::string_view rosetta::ann_json_source<Widget> = R"({
  "title": { "doc": "The widget title" },
  "count": { "doc": "Visible items", "range": [0, 100] },
  "id":    { "readonly": true },
  "mode":  { "combobox": ["fast", "slow"] }
})";
\end{lstlisting}

Same schema, same merge, same effect on every backend. The one rule is that the
specialisation must be visible before \code{rosetta::walk<Widget>()} runs in a
translation unit --- which is exactly why the manifest route bakes it into
\code{bindings.h}.

\subsection{Merge precedence}

Inline annotations come first in the merged pack, JSON after. So a reader that
keeps the last match (\code{ann::get\_or<A>}) lets a JSON entry \emph{override}
an inline one of the same kind, while cumulative readers such as \code{doc} see
both.

A class with no side-car uses the empty primary template and is simply
un-annotated --- no error, no special-casing.

\section{Internal annotations}

\code{rosetta::virtual\_spec} is \textbf{synthesised by the walk}, never written
by users: it marks a method's reflection as virtual or overriding so backends
can emit overridable bindings (pybind11 and N-API trampolines). You only meet it
when writing a custom backend.

\section{Gotchas}

\begin{itemize}
  \item \textbf{Inline annotations make the header C++26-only.} Every
    translation unit that includes it must build with the fork and
    \code{-fannotation-attributes}. Use the side-car when you want off-the-shelf
    builds.
  \item \textbf{Side-car edits need a regeneration}, not just a rebuild.
  \item \code{widget::slider} without a \code{range} falls back to the default
    editor --- a slider needs bounds.
  \item \code{widget::radio} does nothing without a \code{combobox}. It restyles
    the choices; it does not define them.
  \item \code{combobox} is capped at 16 choices. Beyond that a drop-down is the
    wrong widget anyway.
  \item \code{label} changes only the visible text. Scripting backends still
    bind the C++ identifier.
\end{itemize}

\begin{note}
The side-car parser is deliberately small: it takes string values literally (no
escape unescaping) and parses numbers as plain decimals. Enough for the
annotation schema, not a general JSON library.
\end{note}
